\documentclass[12pt, a4paper]{article}

% ---- Packages -----------------------------------------------
\usepackage[margin=1in]{geometry}
\usepackage[doublespacing]{setspace}
\usepackage{amsmath, amssymb}
\usepackage{booktabs}
\usepackage{tabularx}
\usepackage{graphicx}
\usepackage[labelfont=bf, font=small, justification=justified]{caption}
\usepackage[numbers, compress]{natbib}
\usepackage[colorlinks=true, linkcolor=blue, citecolor=blue, urlcolor=blue]{hyperref}
\usepackage{parskip}
\usepackage{lineno}
\usepackage[T1]{fontenc}

\linenumbers
\modulolinenumbers[5]

% Supplementary figure counter
\newcounter{suppfig}
\renewcommand{\thesuppfig}{S\arabic{suppfig}}

% =============================================================
\begin{document}

% ---- Title page ---------------------------------------------
\begin{titlepage}
\begin{center}
{\Large\bfseries The risk was already there: infectious disease
importation at the 2026 FIFA World Cup}

\vspace{1.5cm}

Jose Herrera-Diestra$^{1,*}$

\vspace{0.5cm}

{\small
$^{1}$ Department of [DEPARTMENT], Tarleton State University,
Stephenville, TX 76402, USA\\[4pt]
$^{*}$ Corresponding author.
E-mail: \href{mailto:jherreradiestra@tarleton.edu}{jherreradiestra@tarleton.edu}
}

\vspace{1cm}

\textbf{Running title:} Importation risk at FIFA World Cup 2026

\vspace{0.5cm}

\textbf{Word count (main text):} $\approx$ 3,500

\textbf{Figures:} 6 main + 7 supplementary

\textbf{Tables:} 1
\end{center}
\end{titlepage}

\newpage

% ---- Abstract -----------------------------------------------
\begin{abstract}
\noindent\textbf{Objectives:}
The June 2026 FIFA World Cup group stage will bring an estimated
3.5 million international visitors to 11 US host cities from 48 nations.
We quantified city-level importation risk for five infectious diseases
with explicit uncertainty propagation and a diaspora network extension.

\medskip
\noindent\textbf{Methods:}
Three nested Poisson importation models were constructed: a
pre-tournament baseline (M1), a World Cup-adjusted model (M2), and a
schedule-driven model (M3) routing fans by actual match venues.
Under-reporting correction factors and travel-while-infected
probabilities were sampled via 5,000 Monte Carlo draws. A diaspora
extension used US Census Bureau data to estimate the probability of
imported cases entering co-national social networks at each host city.

\medskip
\noindent\textbf{Results:}
Dengue importation was near-certain ($P(\geq 1) > 0.99$) at Miami,
New York, and Houston under M3 (median $\Lambda = 45.0$, 95\% CI
16.0--92.3). Malaria risk was highest at Atlanta ($P(\geq 1) > 0.99$).
Baseline travel accounted for 86--93\% of expected importations
across all diseases; the World Cup surge added only 7--14\% above
this pre-existing risk. Miami, New York, and Houston showed the
highest probability ($>$0.70) of dengue entering co-national diaspora
networks.

\medskip
\noindent\textbf{Conclusions:}
Dengue is the dominant importation risk, near-certain at Miami, New
York, and Houston; malaria dominates at Atlanta. City-level
surveillance combined with targeted diaspora community outreach is
warranted. The schedule-driven Poisson framework is directly
transferable to future mega-events.

\medskip
\noindent\textbf{Keywords:} importation risk; FIFA World Cup 2026;
dengue; Poisson model; Monte Carlo uncertainty; diaspora networks;
infectious disease surveillance; travel medicine
\end{abstract}

\newpage

% =============================================================
\section{Introduction}
% =============================================================

Mass gathering events concentrate infectious disease importation risk
by drawing large, diverse international populations into a small number
of cities over a short period. The 2026 FIFA World Cup (June 11 --
July 19, 2026) is the first edition with 48 participating nations,
played across 16 venues in the United States, Canada, and Mexico.
The current analysis focuses on the group-stage period (June 2026)
rather than the full tournament window extending into July, as travel
volumes and disease incidence inputs are calibrated to June.
The 11 US host cities --- Atlanta, Boston, Dallas, Houston, Kansas
City, Los Angeles, Miami, New York, Philadelphia, San Francisco, and
Seattle --- are expected to receive approximately 3.5 million international
visitors during the tournament period \citep{NTTO2025}, representing
the largest single sporting event in US history. The geographic
distribution of all 16 WC2026 stadiums across the three host nations
is shown in Supplementary Figure~S4.

Importation risk analyses preceding past FIFA World Cups have taken
different approaches depending on the host country's disease ecology.
Ahead of the 2014 tournament in Brazil, van Panhuis et al.\ quantified
dengue risk for international visitors using travel volumes and local
incidence data \citep{vanPanhuis2014}. For Russia 2018, concerns
centred on respiratory and vaccine-preventable infections in a temperate
host, rather than vector-borne diseases \citep{Pshenichnaya2018}.
Qatar 2022 prompted assessments of multiple importation pathways ---
including dengue, COVID-19, and emerging pathogens --- in a densely
populated host nation receiving visitors from endemic regions
\citep{AlTawfiq2022}. However, those analyses typically relied on
regional-level travel aggregates, reported point estimates without
uncertainty quantification, and did not account for the role of
co-national diaspora communities as secondary amplification networks
after importation.
A concurrent multi-pathogen assessment for WC2026 by \citet{Lessler2026},
using the GLEAM-EPIRisk mobility framework, screened 77 pathogens across
importation risk, outbreak potential, and economic impact, identifying
dengue, malaria, and seasonal respiratory viruses among the priority
pathogens of elevated concern. The present study is complementary to
that breadth-first screening: rather than covering the full pathogen
space, we provide city-specific quantitative importation estimates with
propagated parameter uncertainty and an explicit diaspora seeding pathway
for the five highest-risk diseases.

Three methodological gaps motivate this study. First, country-level
air travel routing data from the Bureau of Transportation Statistics
T-100 programme enable city-specific risk stratification at all 11 US
venues, including cities with historically limited international
surveillance coverage (Atlanta, Kansas City, Seattle). Second,
literature ranges for under-reporting correction factors and
travel-while-infected probabilities span several-fold; propagating
that uncertainty into estimates is essential for honest risk
communication. Third, co-national diaspora communities --- whose
shared language, cultural spaces, and World Cup watch gatherings make
them disproportionately likely to interact with infected travelers
from the same country --- represent a transmission pathway invisible
to standard Poisson arrival models.

We address these gaps with a three-tier nested Poisson importation
model with 5,000-draw Monte Carlo uncertainty for five diseases
(dengue, malaria, measles, pertussis, influenza), coupled to a
diaspora concentration index derived from US Census Bureau American
Community Survey data \citep{Census2026WC, tidycensus}.
The approach is transparent and fully reproducible using only publicly
available data sources, and is designed to be extensible to additional
diseases, geographies, and future mass gathering events.
Its computational simplicity --- requiring only routine air-travel
statistics, published surveillance incidence, and standard statistical
software --- means it can serve as a rapidly deployable first-pass
risk assessment that public health agencies can adapt without
specialized mobility-simulation infrastructure, complementing more
comprehensive frameworks where resources permit.

% =============================================================
\section{Methods}
% =============================================================

We quantified importation risk for five infectious diseases
(dengue, malaria, measles, pertussis, and influenza) at each of the
11 US FIFA World Cup 2026 host cities during June 2026, using three
nested Poisson models that share identical disease data so that
differences in results reflect travel volume assumptions only.

\subsection{Travel data}

Country-of-residence (COR) monthly arrivals were obtained from the
US Customs and Border Protection I-94 programme \citep{CBP_I94}.
City-level routing fractions were derived from BTS T-100 International
Segment (Form 41 Traffic) nonstop passenger records for June
2023--2025 \citep{BTS_T100}, computed as the share of each country's
US-bound passengers that landed at each venue city, averaged over
three years for stability. National Travel and Tourism Office (NTTO)
projections provided country-specific 2026 growth factors for the
12 largest US inbound source markets \citep{NTTO2025}.

\subsection{Three nested model tiers}

\textbf{Baseline (M1)} used COR June 2024 arrivals routed by T-100
fractions with no World Cup growth adjustment, representing a
pre-tournament counterfactual.

\textbf{WC-adjusted (M2)} applied country-specific growth factors
(NTTO projection for top 12 markets; a global World Cup factor of
1.174 for other qualified nations; an empirical background trend
factor for non-qualified countries) to project June 2026 arrivals,
distributed across cities by T-100 fractions.

\textbf{Schedule-driven (M3)} decomposed travel into a World Cup fan
stream --- the marginal increment above background trend, routed
proportionally to each team's actual match venues --- and a
background stream routed by T-100 fractions. This tier correctly
redirects fan travel to Canadian and Mexican venues when teams play
there, producing a more realistic US city risk distribution.
Full mathematical derivations of $N_{c,v}$ for each model tier,
including all intermediate calculations and the formal definition
of the routing fractions, are provided in Appendix~A of the
Supplementary Material.

\subsection{Disease data and Poisson framework}

For each source country $c$, venue city $v$, and disease $d$, the
expected number of infected arrivals is:
\[
  \lambda_{c,v,d}
    = N_{c,v} \;\cdot\; \rho_d \;\cdot\; \frac{C_{c,d}}{P_c}
               \;\cdot\; p_d
\]
where $N_{c,v}$ is projected arrivals, $\rho_d$ is the source-country
under-reporting correction (corrects for undetected cases in surveillance),
$C_{c,d}/P_c$ is reported per-capita incidence ($P_c$ are 2026
national population estimates \citep{WorldBank2024pop}), and $p_d$ is the
probability that an infected individual presents with mild or no
symptoms and therefore remains travel-capable (severely symptomatic
individuals are assumed to self-select out of international travel). City-level intensity
and importation probability follow as
$\Lambda_{v,d} = \sum_c \lambda_{c,v,d}$ and
$P(\geq 1)_{v,d} = 1 - e^{-\Lambda_{v,d}}$.

Disease-specific incidence was sourced from: WHO Global Dengue
Surveillance dataset (monthly; June 2024--2025 average)
\citep{WHO_dengue}; WHO World Malaria Report 2024 (annual incidence
per 1,000 population) \citep{whoWMR2024}; WHO Immunization Data
(annual incidence per million for measles and pertussis)
\citep{WHO_immunization}; WHO FluNet GISRS (weekly positive specimens,
ISO weeks 22--26, 2023--2025 average). Literature-derived disease
parameters are summarised in Table~\ref{tab:params}.

\begin{table}[htbp]
\centering
\caption{Disease parameters and Monte Carlo uncertainty bounds used
  in all three model tiers. For each of 5,000 iterations,
  $\rho_d \sim \mathrm{Uniform}(\rho_{\min},\rho_{\max})$ and
  $p_d \sim \mathrm{Uniform}(p_{\min},p_{\max})$ are drawn
  independently. All reported estimates are Monte Carlo medians
  with 95\% uncertainty intervals. Dengue and influenza incidence
  are June-month values; malaria, measles, and pertussis use
  annual rates divided by 12 to approximate June incidence.}
\label{tab:params}
\small
\setlength{\tabcolsep}{6pt}
\begin{tabular}{l rr rr l}
\toprule
&
  \multicolumn{2}{c}{$\boldsymbol{\rho_d}$} &
  \multicolumn{2}{c}{$\boldsymbol{p_d}$} & \\
\cmidrule(lr){2-3}\cmidrule(lr){4-5}
\textbf{Disease} &
  \textbf{min} & \textbf{max} &
  \textbf{min} & \textbf{max} &
  \textbf{References} \\
\midrule
Dengue    & 0.06 & 0.26 & 0.30 & 0.70 & \citep{bhatt2013,undurraga2013} \\[3pt]
Malaria   & 0.10 & 0.35 & 0.10 & 0.50 & \citep{whoWMR2024,lindblade2013,topazian2020} \\[3pt]
Measles   & 0.40 & 0.80 & 0.02 & 0.10 & \citep{simons2012,whittle1999,helfand1998} \\[3pt]
Pertussis & 0.01 & 0.10 & 0.50 & 0.90 & \citep{mclaughlin2016,craig2020,althouse2015} \\[3pt]
Influenza & 0.01 & 0.10 & 0.30 & 0.70 & \citep{mccarthy2020,hayward2014,carrat2008,leung2015} \\
\bottomrule
\end{tabular}
\end{table}

\subsection{Monte Carlo uncertainty quantification}

The under-reporting correction $\rho_d$ and travel-while-infected
probability $p_d$ each have wide literature ranges reflecting
heterogeneity across surveillance systems and epidemiological settings.
We treated both as independent Uniform random variables within their
literature bounds (Table~\ref{tab:params}) and drew 5,000 samples per
disease. For each draw, all city-level $\Lambda_{v,d}$ values were
recomputed. Reported estimates are Monte Carlo medians with 95\%
uncertainty intervals (2.5th--97.5th percentiles). Because $\rho_d$
and $p_d$ enter multiplicatively, uncertainty intervals are asymmetric:
diseases whose reference $\rho$ lies near the lower literature bound
(dengue) have strongly right-skewed intervals, indicating the reported
estimate is conservative; diseases at the upper bound (pertussis) have
left-skewed intervals (Supplementary Figure~S3).

\subsection{Diaspora network extension}

Standard Poisson models treat all arrivals identically regardless of
the social environment they enter. We extended the framework to
estimate the probability that imported cases reach co-national
diaspora networks at each host city.

Diaspora concentration was computed from US Census Bureau American
Community Survey 5-year estimates (2019--2023), Table B05006, via the
\texttt{tidycensus} R package \citep{tidycensus,Census2026WC},
defined as the share of each metro area's total foreign-born
population born in each source country.

\textbf{Mechanism A (local social mixing).} Each infected arrival
independently contacts the co-national community with a probability
proportional to the diaspora concentration index. This yields a
city-level probability that at least one imported case enters a
co-national network, directly comparable to the main model's
$P(\geq 1)$ metric.

\textbf{Mechanism B (return seeding).} US diaspora members may
travel domestically to match venues and carry infection back to
their home city. A secondary seeding index ranks hub cities by
their exposure potential, combining importation intensity at the
match venue with the diaspora concentration in the hub city.

Full mathematical derivations of both mechanisms, including the
Poisson thinning argument, all intermediate quantities, and the
formal relationship between the diaspora-adjusted and standard
importation probabilities, are provided in Appendix~B of the
Supplementary Material.

All analyses were conducted in R (version 4.4) using packages
\texttt{tidyverse}, \texttt{readxl}, \texttt{cowplot}, and
\texttt{tidycensus}.

% =============================================================
\section{Results}
% =============================================================

Unless otherwise noted, results below refer to the schedule-driven
model (M3), which allocates World Cup fan travel by match schedule
and is the most faithful representation of actual visitor flows; a
direct comparison across all three model tiers is provided in
Section~\ref{sec:comparison}.

\subsection{Importation risk by disease and host city}

Under the schedule-driven model (M3), dengue importation was
near-certain ($P(\geq 1) > 0.99$) at Miami, New York, and Houston
(Figure~\ref{fig:heatmap}; equivalent heatmaps for the Baseline (M1)
and WC-adjusted (M2) models are shown in Supplementary Figures~S1
and S2). Miami had the highest absolute
importation intensity ($\Lambda = 15.8$, 95\% CI 5.6--32.3),
followed by New York (10.5, 95\% CI 3.7--21.5) and Houston (5.6,
95\% CI 2.0--11.4). Dengue importation probabilities exceeded 0.95 at
Dallas and Atlanta, while smaller venues --- Boston, Seattle, and
Kansas City --- showed substantially lower probabilities. The
11-city dengue total under M3 was 45.0 (95\% CI 16.0--92.3).
Malaria risk was highest at Atlanta ($P(\geq 1) > 0.99$),
reflecting the concentration of Sub-Saharan African air routes
through Hartsfield-Jackson airport; the 11-city malaria total was
16.4 (95\% CI 4.9--39.1). Pertussis and influenza
showed intermediate probabilities. Measles had near-zero expected
importation intensity ($\Lambda = 0.05$, 95\% CI 0.02--0.10) because
WHO-reported per-capita incidence is very low across most World Cup
source countries; US measles cases arise primarily from outbreak
chains seeded by a small number of imports rather than the independent
Poisson arrivals estimated here.

\begin{figure}[htbp]
  \centering
  \includegraphics[width=\textwidth]{fig1_risk_heatmap_IJID.png}
  \caption{\textbf{Importation probability for five infectious diseases
    at 11 US FIFA World Cup 2026 host cities.}
    Cells show $P(\geq 1)$, the probability of at least one imported
    case arriving at each city in June 2026, under the schedule-driven
    model (M3) Monte Carlo median. Cities are ordered left-to-right by
    total importation intensity across all five diseases. Dengue
    $P(\geq 1) > 0.99$ at Miami, New York, and Houston; malaria exceeds
    0.99 at Atlanta.
    Values $<$0.01 shown as $<$0.01; values $>$0.99 as $>$0.99.}
  \label{fig:heatmap}
\end{figure}

\subsection{Expected importation intensity and parameter uncertainty}

Figure~\ref{fig:lambdaci} shows Monte Carlo median $\Lambda$ and 95\%
uncertainty intervals for each disease and city. Dengue intensities
are one to two orders of magnitude larger than those of other diseases.
Interval asymmetry is disease-specific: for dengue, the upper interval
bound is roughly 2$\times$ the median while the lower bound is
$\approx$0.35$\times$ the median, indicating that the central estimate
is near the lower plausible value. For pertussis, the interval
collapses toward the lower end, confirming that estimates are
conservative upper bounds (the reference detection rate equals the
upper literature limit).

\begin{figure}[htbp]
  \centering
  \includegraphics[width=\textwidth]{fig2_lambda_ci_IJID.png}
  \caption{\textbf{Expected importation intensity with 95\% Monte Carlo
    uncertainty intervals, by disease and host city (schedule-driven
    model, M3).}
    Points are Monte Carlo medians; horizontal bars are 2.5th--97.5th
    percentiles from 5,000 Uniform draws on $\rho_d$ and $p_d$.
    Each panel has a free x-axis; note the order-of-magnitude scale
    difference between dengue and other diseases. Asymmetric intervals
    reflect whether the reference parameter value lies near the lower
    (dengue, right-skewed) or upper (pertussis, left-skewed) literature
    bound (Supplementary Figure~S3).}
  \label{fig:lambdaci}
\end{figure}

\subsection{Effect of the World Cup travel surge across model tiers}
\label{sec:comparison}

The World Cup travel surge (M1 $\to$ M2) increased the 11-city dengue
total by approximately 17\% (39.7 to 46.6; Figure~\ref{fig:comparison}).
The schedule-driven model (M3: 45.0) fell slightly below the
WC-adjusted model (M2: 46.6) because the fan increment for teams
playing in Canada and Mexico was redirected to non-US venues,
reducing US city totals. This difference --- absent from M2, which
allocates all World Cup growth to US cities via T-100 routing ---
is structurally important: M3 is the more accurate estimate of the
US-specific burden. Comparable patterns held for malaria (M1: 14.7,
M2: 16.6, M3: 16.4) and other diseases.

\begin{figure}[htbp]
  \centering
  \includegraphics[width=\textwidth]{fig3_model_comparison_IJID.png}
  \caption{\textbf{Three-model comparison of expected importation
    intensity for the six highest-risk US host cities.}
    Bars show Monte Carlo median $\Lambda$; error bars show 95\%
    uncertainty intervals. M1 = Baseline (COR June 2024, no World Cup
    adjustment); M2 = WC-adjusted (2026 projections, T-100 routing for
    all travellers); M3 = Schedule-driven (World Cup fans routed by
    match schedule, background travellers by T-100). Panels use free
    y-axes; note differences in scale across diseases.}
  \label{fig:comparison}
\end{figure}

Figure~\ref{fig:excess} decomposes the 11-city total into baseline
travel (M1) and the World Cup-specific increment (M3 $-$ M1) for
each disease. Baseline travel accounts for 86--93\% of total expected
importations across all diseases, with the WC increment contributing
7--14\% above the pre-existing risk. This finding is
epidemiologically important: the dominant importation burden exists
independently of the World Cup, and surveillance systems calibrated
only to the tournament window will miss the majority of the risk.

\begin{figure}[htbp]
  \centering
  \includegraphics[width=\textwidth]{fig4_wc_excess_IJID.png}
  \caption{\textbf{World Cup travel increment above baseline
    importation by disease.}
    Each panel shows the 11-city Monte Carlo median $\Lambda$ split
    into baseline travel (M1, blue) and the schedule-driven WC
    increment (M3 $-$ M1, red). Panels use free y-axes. Percentage
    labels indicate the WC increment as a share of the M1 baseline.
    Baseline travel dominates across all five diseases, contributing
    86--93\% of total expected importations.}
  \label{fig:excess}
\end{figure}

\subsection{Source-country contributions}

Brazil and Colombia were the dominant dengue contributors across New
York, Miami, and Los Angeles (Figure~\ref{fig:countries}). Nigeria
and the Democratic Republic of Congo drove malaria risk at Atlanta
and Houston. For dengue at World Cup fan countries, contributions
were split between a fan-specific stream (concentrated at match
venues) and a background stream (distributed by routine T-100 routing),
confirming that the surge amplification is real but built on
pre-existing high baseline travel from endemic source countries.
Supplementary Figure~S6 ranks the same top-10 source countries by
diaspora-weighted importation intensity ($\Omega^A$) rather than raw
$\lambda$; countries with large co-national communities concentrated
in their primary destination cities rank higher in $\Omega^A$ than
in $\lambda$, revealing an additional layer of community-level
exposure not visible from travel volumes alone.

\begin{figure}[htbp]
  \centering
  \includegraphics[width=\textwidth]{fig5_country_drivers_IJID.png}
  \caption{\textbf{Top 10 source-country contributors to expected
    importation intensity under the schedule-driven model (M3),
    by disease.}
    Bars show expected importations $\lambda_{c,\cdot,d}$ summed
    across all 11 US venue cities; error bars show 95\% Monte Carlo
    uncertainty intervals. Colors denote world region.
    (A)~Dengue: Latin America dominates, led by Brazil and Colombia.
    (B)~Influenza: contributions broadly distributed across source
    countries, consistent with global co-circulation during June.
    (C)~Pertussis: Europe and Latin America lead, reflecting
    waning vaccine-induced immunity in adult traveler populations.
    (D)~Malaria: Sub-Saharan Africa dominates, led by Nigeria and
    the Democratic Republic of Congo.
    (E)~Measles: near-zero contributions across all source countries,
    consistent with the very low global per-capita incidence used
    in the model.}
  \label{fig:countries}
\end{figure}

\subsection{Diaspora network seeding risk}

Figure~\ref{fig:diaspora} summarises both diaspora risk mechanisms
for dengue and malaria, the two diseases with meaningful diaspora
signal. Panel A shows $P^A(\geq 1)_{v,d}$, the probability that at
least one imported case enters a co-national diaspora social network
at each host city. For dengue, Miami, New York, and Houston exceeded
$P^A(\geq 1) = 0.70$, driven respectively by large Haitian and
Brazilian communities in Miami ($\kappa = 0.088$ and $0.079$), large
Ecuadorian and Colombian communities in New York, and large Venezuelan
and Colombian communities in Houston. The gap between
$P(\geq 1) \approx 1$ (near-certain arrival) and
$P^A(\geq 1) \approx 0.60\text{--}0.80$ (diaspora network entry)
quantifies risk dispersed into the general population rather than
concentrated in a co-national community. For malaria, Atlanta showed
the highest $P^A(\geq 1)$, consistent with its large Nigerian and
Ghanaian diaspora populations relative to the total foreign-born
community.

Panel B shows the Mechanism B secondary seeding index ($\Omega^B$)
for hub cities outside the WC venue network. Cities not hosting
matches but with large diaspora populations from high-dengue or
high-malaria source nations face non-trivial secondary exposure
through returning co-national fans. Results for all five diseases
are provided in Supplementary Figure~S5. The full Mechanism A
diaspora probability ($P^A(\geq 1)$) across all five diseases and
11 host cities is provided in Supplementary Figure~S7.

\begin{figure}[htbp]
  \centering
  \includegraphics[width=\textwidth]{fig6_diaspora_combined_IJID.png}
  \caption{\textbf{Diaspora network seeding risk for dengue and
    malaria.}
    \textbf{Panel A (Mechanism A):} $P^A(\geq 1)_{v,d} =
    1 - e^{-\Omega^A_{v,d}}$, the probability that at least one
    imported case enters a co-national diaspora network at each host
    city. Derived from $\Omega^A_{v,d} = \sum_c
    \lambda_{c,v,d}\,\kappa_{c,v}$, where
    $\kappa_{c,v} = D_{c,v}/\mathrm{FB}_v$ (ACS 2019--2023).
    $P^A(\geq 1)$ is always $\leq P(\geq 1)$
    (Figure~\ref{fig:heatmap}); the gap represents risk entering the
    general population.
    \textbf{Panel B (Mechanism B):} Secondary seeding index
    $\Omega^B$ for the top hub cities per disease. Hub cities are
    non-venue cities where diaspora members may attend games and
    return home, potentially seeding local transmission.}
  \label{fig:diaspora}
\end{figure}

% =============================================================
\section{Discussion}
% =============================================================

Dengue importation is near-certain ($P(\geq 1) > 0.99$) at Miami,
New York, and Houston, and exceeds $P(\geq 1) = 0.95$ at Dallas
and Atlanta, driven by high reported burden in WC-qualified Latin
American nations during June, large and growing travel volumes, and
a relatively long travel-while-infected window. Critically, the
right-skewed uncertainty interval for dengue --- where the reference
detection rate lies near the lower literature bound --- means the
central estimate is likely conservative. The true 11-city June total
could plausibly exceed 90 infected arrivals under the upper tail of
the parameter distribution.

These estimates are consistent with the record travel-associated
dengue burden observed in the US during 2024 (over 3,798 cases
nationally, a 359\% increase over the historical average). Although
2025 incidence moderated, the large and growing US diaspora from
endemic Latin American countries sustains a high baseline risk that
the World Cup travel surge amplifies further.

The comparison across model tiers reveals a structural insight
not available from standard event-based risk assessments: M3 is
strictly more accurate for the 11 US cities than M2 because it
correctly routes some World Cup fan travel to non-US venues in
Canada and Mexico. The M2-minus-M3 gap is modest for dengue
($\approx$3\%) but represents a real redistribution that matters
for city-specific public health preparedness. Surveillance and
response resources calibrated to M3 estimates will be better
matched to the actual risk at US venues.

Perhaps the most policy-relevant finding is that baseline travel
already accounts for 86--93\% of total expected importations across
all five diseases (Figure~\ref{fig:excess}). The World Cup adds
7--14\% above a pre-existing risk that is present every June
regardless of the tournament. This implies that targeted
tournament-period surveillance, while valuable, cannot substitute
for year-round monitoring of international travelers arriving at US
airports from endemic regions. The marginal contribution of the WC
travel surge is real but modest relative to the structural baseline
risk --- a distinction that is obscured by event-specific assessments
that do not include a pre-tournament counterfactual.

The diaspora extension introduces a distinction absent from
standard importation models: whether an imported case enters a
concentrated co-national social network or disperses into the
general population. This distinction is operationally important
because co-national networks have higher contact rates for disease
transmission, may face language and access barriers to healthcare,
and are often the focus of World Cup watch-party gatherings that
could amplify early transmission chains. The $P^A(\geq 1)$ metric
provides a directly actionable quantity for targeted community
health outreach that complements city-level $P(\geq 1)$ monitoring.

Several limitations should be noted. First, the Poisson independence
assumption may understate clustering risk from coordinated group travel
by fans from the same country. Second, $\rho_d$ corrects for
under-reporting in source-country surveillance but does not model
US detection rates; our $\Lambda$ estimates represent true infected
arrivals, not detected US cases, so comparisons with CDC surveillance
tallies require additional adjustment. Third, $C_{c,d}/P_c$ uses
total-population incidence without age stratification, which introduces
directional biases that vary by disease: malaria risk from high-endemic
sub-Saharan African settings is likely overestimated because pediatric
cases dominate reported burden while adult travelers typically carry
partial acquired immunity through repeated prior exposure, whereas for
pertussis the bias is smaller because waning vaccine-induced immunity
has progressively shifted burden toward the young-adult traveler
demographic. Fourth, diaspora concentration $\kappa_{c,v}$ is derived
from ACS 2019--2023 data spanning the COVID-19 pandemic, which altered
immigration patterns; values may not perfectly reflect the June 2026
population distribution. Fifth, the analysis covers only the 11 US
host cities; Canadian and Mexican venues, which receive meaningful
World Cup fan flows, are outside the primary scope. Sixth, the near-zero measles estimate reflects
the model's scope --- it estimates independent seeding events, not
outbreak chains; US measles case counts typically arise from a small
number of imports followed by local transmission in undervaccinated
communities, a process not captured by the Poisson framework.
Finally, because $p_d$ enters multiplicatively it is a key driver
of uncertainty: a one-parameter sensitivity analysis (Table~S1)
confirms that malaria and measles estimates are most sensitive to
this parameter (fold-change up to 5.0$\times$ across the literature
range), while pertussis is least sensitive (1.8$\times$); absolute
importation counts at each $p_d$ level are provided in the
analysis code repository.

% =============================================================
\section{Conclusions}
% =============================================================

Dengue importation is near-certain at Miami, New York, and Houston,
and represents a substantial risk at Dallas and Atlanta; malaria is
the dominant threat at Atlanta. Baseline travel accounts for
86--93\% of total expected importations across all diseases,
confirming that the dominant risk exists year-round independently of
the tournament, and that the World Cup travel surge contributes a
real but modest 7--14\% increment above that pre-existing burden.
A two-tier surveillance strategy is warranted: heightened city-level
monitoring for dengue and malaria combined with targeted community
outreach within co-national diaspora networks at Miami, New York,
and Houston, where the probability of a case entering a co-national
social network exceeds 0.70. The schedule-driven Poisson framework
is directly reusable for future mega-events by substituting updated
travel volumes and disease incidence data, and its reliance on
publicly available inputs makes it a practical first-pass tool for
public health agencies planning risk assessments for any large
international gathering.

% =============================================================
% DECLARATIONS
% =============================================================
\section*{Conflict of interest}
The author declares no conflict of interest.

\section*{Funding}
[Add funding information here.]

\section*{Data availability}
All disease data are publicly available from WHO repositories cited
below. BTS T-100 and COR/I-94 travel data are available from the
Bureau of Transportation Statistics and US Customs and Border
Protection respectively. Diaspora data were obtained from the
US Census Bureau ACS 5-year estimates (Table B05006).
Analysis code is available at \url{https://github.com/jdiestra1977/FIFA\_worldCup\_2026\_risk}.

% =============================================================
% REFERENCES
% =============================================================
\bibliographystyle{vancouver}
\bibliography{references_IJID}

% =============================================================
% SUPPLEMENTARY MATERIAL
% =============================================================
\newpage

\begin{center}
{\large\bfseries Supplementary Material}
\end{center}

\setcounter{suppfig}{0}

% ---- Appendix A ---------------------------------------------
\subsection*{Appendix A. Mathematical specification of travel flow
estimates $N_{c,v}$}

\subsubsection*{A.1 Notation}

\begin{center}
\small
\begin{tabular}{p{2.2cm} p{10cm}}
\toprule
\textbf{Symbol} & \textbf{Definition} \\
\midrule
$c$ & Source country ($c = 1,\ldots,C$) \\
$v$ & US venue city ($v = 1,\ldots,11$) \\
$u$ & Any WC venue (US, Canadian, or Mexican) \\
$A_c$ & June 2024 COR arrivals from country $c$ to the US
        (I-94 programme, CBP) \\
$T_{c,v}$ & Mean June passengers from country $c$ to US city $v$,
        averaged over June 2023--2025 (BTS T-100 Form 41 data) \\
$f_{c,v}^{T}$ & T-100 routing fraction: share of country $c$'s
        US-bound passengers landing at city $v$ (see A.2).
        By construction $\sum_{v} f_{c,v}^{T} = 1$. \\
$\phi_c$ & Country-specific 2026 growth factor (see A.3) \\
$g_c$ & Background (non-WC) trend component of $\phi_c$,
        representing non-WC travel growth from 2024 to 2026 \\
$m_{c,v}$ & Number of group-stage matches that country $c$
        plays at US venue $v$ \\
$M_c$ & Total group-stage matches for country $c$
        ($M_c = \sum_u m_{c,u} = 3$ for all qualified teams) \\
$\omega_{c,v}$ & Schedule-driven routing fraction for WC fans
        from country $c$ to US venue $v$ (see A.4) \\
$N_{c,v}$ & Projected June 2026 arrivals from country $c$
        to venue city $v$ (model-tier specific) \\
\bottomrule
\end{tabular}
\end{center}

\subsubsection*{A.2 T-100 routing fraction $f_{c,v}^{T}$}

The T-100 routing fraction gives the proportion of country $c$'s
US-bound passengers that land at each venue city $v$:
\[
  f_{c,v}^{T} = \frac{T_{c,v}}{\displaystyle\sum_{v'=1}^{11} T_{c,v'}}
\]
where $T_{c,v}$ is the mean number of June passengers from country
$c$ to US city $v$ over 2023--2025 (BTS T-100 Form 41 data). By
construction, $\sum_{v} f_{c,v}^{T} = 1$ for every country $c$.
These fractions are held fixed across all three model tiers; only
the total volume $A_c \cdot \phi_c$ changes between M1, M2, and M3.

\subsubsection*{A.3 Growth factor $\phi_c$ and its components}

The growth factor
\[
  \phi_c = \frac{V_{c,2026}}{V_{c,2024}}
\]
scales June 2024 baseline arrivals to projected June 2026 levels,
where $V_{c,t}$ denotes projected total June visitors from country
$c$ in year $t$. Three country classes are distinguished:

\begin{enumerate}
  \item \textbf{Top-12 inbound markets} (Australia, Brazil, Canada,
    China, France, Germany, India, Italy, Japan, Mexico, South Korea,
    UK) --- NTTO country-specific projections:
    \[
      \phi_c = \frac{A_{c,2026}^{\text{NTTO}}}{A_{c,2024}}
    \]
    The World Cup effect is already embedded in the NTTO forecast.
    The background (non-WC) trend component $g_c$ is derived from
    the pre-tournament NTTO trend excluding the tournament increment.

  \item \textbf{Other WC-qualified nations} --- a uniform World Cup
    growth factor is applied:
    \[
      \phi_c = \phi^{\text{WC}} = \frac{85{,}017}{72{,}390} = 1.174
    \]
    reflecting the aggregate fan-travel surge for smaller qualified
    markets (NTTO 2025 report). The background trend component is
    $g_c = \phi^{\text{bg}} = 1.077$, so the WC-specific fan
    increment is $\phi_c - g_c = 0.097$.

  \item \textbf{Non-qualified nations} --- no WC adjustment is made:
    \[
      \phi_c = g_c = \phi^{\text{bg}} = 1.077
    \]
    where $\phi^{\text{bg}}$ is the data-driven background growth
    rate estimated from COR and T-100 data (June 2023--2025).
    By construction, $\phi_c - g_c = 0$ for non-qualified countries.
\end{enumerate}

\subsubsection*{A.4 Schedule-driven routing fraction $\omega_{c,v}$}

For WC-qualified nation $c$, the fan stream is routed proportionally
to the share of its group-stage matches played at each US venue:
\[
  \omega_{c,v} = \frac{m_{c,v}}{M_c} = \frac{m_{c,v}}{3}
  \quad \text{(US venues only)}
\]
If country $c$ plays no matches at US venue $v$ then
$\omega_{c,v} = 0$. Matches played at Canadian or Mexican venues
contribute to $M_c$ but not to any US $\omega_{c,v}$, so:
\[
  \sum_{v=1}^{11} \omega_{c,v}
    = \frac{\text{(US matches of }c)}{3} \;\leq\; 1
\]
with strict inequality for nations playing one or more matches
outside the US. For non-qualified nations no WC fan stream exists
and $\omega_{c,v}$ is not used.

\subsubsection*{A.5 Travel decomposition for M3}

Model 3 separates total projected 2026 arrivals from country $c$
into a background stream and a WC fan stream:
\[
  N_c^{\text{WC}} = A_c \cdot \max\!\bigl(0,\;\phi_c - g_c\bigr)
  \qquad
  N_c^{\text{bg}} = A_c \cdot g_c
\]
Note that $N_c^{\text{WC}} + N_c^{\text{bg}} = A_c \cdot \phi_c$,
so the decomposition is exact and exhaustive. For non-qualified
countries $\phi_c = g_c$, giving $N_c^{\text{WC}} = 0$. For
typical WC-qualified nations the fan increment represents
approximately $(1.174 - 1.077)/1.174 \approx 8.3\%$ of total
projected 2026 arrivals.

\subsubsection*{A.6 Travel flow $N_{c,v}$ by model tier}

\paragraph{M1 --- Baseline.}
Raw June 2024 arrivals distributed by T-100 routing fractions;
no growth factor applied.
\[
  \boxed{N_{c,v}^{M1} = A_c \cdot f_{c,v}^{T}}
\]

\paragraph{M2 --- WC-adjusted.}
Projected 2026 arrivals (incorporating the full growth factor
$\phi_c$) distributed by T-100 routing fractions. All WC fan
travel is assigned to US cities regardless of whether matches
are played in Canada or Mexico.
\[
  \boxed{N_{c,v}^{M2} = A_c \cdot \phi_c \cdot f_{c,v}^{T}}
\]

\paragraph{M3 --- Schedule-driven.}
The background stream is T-100 routed; the WC fan stream is
routed by the match schedule:
\[
  \boxed{N_{c,v}^{M3}
    = \underbrace{N_c^{\text{bg}} \cdot f_{c,v}^{T}}_{\text{background stream}}
      \;+\;
      \underbrace{N_c^{\text{WC}} \cdot \omega_{c,v}}_{\text{WC fan stream}}
    = A_c \Bigl[
        g_c \cdot f_{c,v}^{T}
        + (\phi_c - g_c) \cdot \omega_{c,v}
      \Bigr]}
\]
For non-qualified nations $\phi_c = g_c$, so the fan stream
vanishes and $N_{c,v}^{M3} = A_c \cdot g_c \cdot f_{c,v}^{T}$.

\subsubsection*{A.7 Why $N_{c,v}^{M3} \leq N_{c,v}^{M2}$ for
US cities}

M2 routes the entire WC increment $(\phi_c - g_c)$ through T-100
fractions, which sum to 1 over US cities. M3 routes the same
increment through schedule fractions, which sum to $\leq 1$ over
US cities (strictly $<1$ whenever country $c$ plays matches in
Canada or Mexico). Therefore:
\[
  \sum_{v} N_{c,v}^{M3}
    = A_c \Bigl[
        g_c \underbrace{\sum_v f_{c,v}^T}_{=1}
        + (\phi_c - g_c)
          \underbrace{\sum_v \omega_{c,v}}_{\leq 1}
      \Bigr]
  \;\leq\;
  A_c \cdot \phi_c
    = \sum_v N_{c,v}^{M2}
\]
The gap equals the fan travel redirected to non-US WC venues
(Canada, Mexico), which is why the M3 US-city total lies below
M2 for all qualified nations with at least one non-US match.
Importation probability heatmaps for M1 and M2 are shown in
Supplementary Figures~S1 and S2; these complement the M3 heatmap
in Figure~\ref{fig:heatmap} and the three-model comparison in
Figure~\ref{fig:comparison}.

\newpage

% ---- Appendix B ---------------------------------------------
\subsection*{Appendix B. Mathematical specification of the diaspora
network extension}

\subsubsection*{B.1 Notation}

\begin{center}
\small
\begin{tabular}{p{2.8cm} p{9.5cm}}
\toprule
\textbf{Symbol} & \textbf{Definition} \\
\midrule
$D_{c,v}$            & Number of US residents born in country $c$
                       living in metro area $v$
                       (ACS 5-year estimates 2019--2023, Table B05006) \\
$\mathrm{FB}_{v}$    & Total foreign-born population of metro area $v$
                       (row total of Table B05006) \\
$\kappa_{c,v}$       & Diaspora concentration index: share of metro
                       $v$'s foreign-born population born in country $c$ \\
$\Omega^A_{c,v,d}$   & Expected number of infected arrivals from $c$
                       at city $v$ that enter the co-national network
                       (Mechanism A) \\
$P^A(\geq 1)_{v,d}$  & Probability that at least one infected arrival
                       enters any co-national network in city $v$ \\
$v_{\text{match}}$   & WC venue city where country $c$ plays \\
$v_{\text{home}}$    & Diaspora hub city (may differ from
                       $v_{\text{match}}$) \\
$\Omega^B_{c,v_m,v_h,d}$ & Secondary seeding index for hub city
                       $v_h$ arising from match at $v_m$ (Mechanism B) \\
\bottomrule
\end{tabular}
\end{center}

\subsubsection*{B.2 Diaspora concentration index $\kappa_{c,v}$}

For each source country $c$ and venue metro area $v$:
\[
  \kappa_{c,v} = \frac{D_{c,v}}{\mathrm{FB}_{v}}
\]
$\kappa_{c,v} \in [0,1]$ measures how concentrated the country-$c$
community is \emph{within} the immigrant population of city $v$ ---
the relevant social network for co-national mixing. By construction,
$\sum_c \kappa_{c,v} = 1$ for each $v$.

\subsubsection*{B.3 Mechanism A: local social mixing}

Let $X_{c,v,d} \sim \text{Poisson}(\lambda_{c,v,d})$ be the number
of infected arrivals from country $c$ at city $v$. Under the
proportional-mixing assumption, each arrival independently enters
the co-national network with probability $\kappa_{c,v}$. By the
Poisson thinning theorem, the thinned count is:
\[
  Y_{c,v,d} \sim \text{Poisson}(\lambda_{c,v,d} \cdot \kappa_{c,v})
\]
We define the community-weighted importation intensity:
\[
  \Omega^A_{c,v,d} = \lambda_{c,v,d} \cdot \kappa_{c,v}
    = \lambda_{c,v,d} \cdot \frac{D_{c,v}}{\mathrm{FB}_{v}}
\]
Because $Y_{c,v,d}$ values are independent Poisson random variables,
the city-level aggregate $\Omega^A_{v,d} = \sum_c \Omega^A_{c,v,d}$
is also a valid Poisson rate. The probability that at least one
infected arrival enters any co-national network in city $v$ is:
\[
  P^A(\geq 1)_{v,d} = 1 - e^{-\Omega^A_{v,d}}
\]
This is directly comparable to the main-model probability
$P(\geq 1)_{v,d} = 1 - e^{-\Lambda_{v,d}}$. Since
$\kappa_{c,v} \leq 1$ for all $c,v$, it follows that
$\Omega^A_{v,d} \leq \Lambda_{v,d}$ and therefore
$P^A(\geq 1)_{v,d} \leq P(\geq 1)_{v,d}$ always: the
diaspora-network probability is a lower bound on the arrival
probability, never exceeding it.

\subsubsection*{B.4 Mechanism B: diaspora return seeding}

US-resident diaspora members from country $c$ living in hub city
$v_{\text{home}}$ may travel domestically to $v_{\text{match}}$
(where country $c$ plays), be exposed to infected international
arrivals, and seed transmission upon return. The secondary seeding
index is:
\[
  \Omega^B_{c,\,v_{\text{match}},\,v_{\text{home}},\,d}
    = \lambda_{c,\,v_{\text{match}},\,d}
      \cdot \kappa_{c,\,v_{\text{home}}}
\]
Note that in Mechanism A both $\lambda$ and $\kappa$ reference the
\emph{same} city; in Mechanism B they reference \emph{different}
cities ($v_{\text{match}}$ and $v_{\text{home}}$, respectively).
$\Omega^B$ is therefore a \textbf{relative seeding index} that
ranks hub cities by exposure potential rather than a strict Poisson
arrival rate. The aggregate index for hub city $v_{\text{home}}$
across all source countries and match venues is:
\[
  \Omega^B_{v_{\text{home}},\,d}
    = \sum_c \sum_{v_{\text{match}} \neq v_{\text{home}}}
      \Omega^B_{c,\,v_{\text{match}},\,v_{\text{home}},\,d}
\]

\newpage

% ---- Table S1 -----------------------------------------------
\subsection*{Table S1. Sensitivity of expected importations to the
mild-symptom travel probability $p_d$.}

Because $p_d$ enters the model multiplicatively
($\lambda_{c,v,d} = N_{c,v}\cdot\rho_d\cdot C_{c,d}/P_c\cdot p_d$),
fixing $\rho_d$ at its literature midpoint and varying $p_d$ produces
importation totals that scale exactly as $p_d/p_{\text{mid}}$. The
table below reports these relative values ($\Lambda_{\text{rel}}$,
normalised to 1.00 at $p_{\text{mid}}$) and the exact fold-change
$p_{\max}/p_{\min}$ for all five modelled diseases. Absolute 11-city
importation counts at each $p_d$ level are available in the
analysis code at \url{https://github.com/jdiestra1977/FIFA\_worldCup\_2026\_risk}.

\medskip
\begin{center}
\small
\begin{tabular}{lcccrrrr}
\toprule
\textbf{Disease} &
  $\boldsymbol{p_{\min}}$ &
  $\boldsymbol{p_{\text{mid}}}$ &
  $\boldsymbol{p_{\max}}$ &
  $\boldsymbol{\Lambda(p_{\min})}$ &
  $\boldsymbol{\Lambda(p_{\text{mid}})}$ &
  $\boldsymbol{\Lambda(p_{\max})}$ &
  \textbf{Fold} \\
\midrule
Dengue    & 0.30 & 0.50 & 0.70 & 29.70 & 49.50 & 69.29 & 2.3$\times$ \\
Malaria   & 0.10 & 0.30 & 0.50 &  6.08 & 18.25 & 30.42 & 5.0$\times$ \\
Measles   & 0.02 & 0.06 & 0.10 &  0.02 &  0.05 &  0.08 & 5.0$\times$ \\
Pertussis & 0.50 & 0.70 & 0.90 &  1.49 &  2.08 &  2.68 & 1.8$\times$ \\
Influenza & 0.30 & 0.50 & 0.70 &  1.19 &  1.98 &  2.78 & 2.3$\times$ \\
\bottomrule
\end{tabular}
\end{center}

\medskip
\noindent
$\Lambda$ = 11-city total expected importations (M3 schedule-driven
model), computed deterministically with $\rho_d$ fixed at its
literature midpoint $({\rho_{\min}+\rho_{\max}})/2$.
Fold $= \Lambda(p_{\max})/\Lambda(p_{\min}) = p_{\max}/p_{\min}$ (exact).
Malaria and measles carry the largest sensitivity to $p_d$; pertussis
the smallest.

\newpage

% ---- Figure S1 ----------------------------------------------
\refstepcounter{suppfig}
\begin{figure}[htbp]
  \centering
  \includegraphics[width=\textwidth]{figS1_heatmap_baseline_IJID.png}
  \caption*{\textbf{Figure S\arabic{suppfig}.}
    \textbf{Importation probability under the Baseline model (M1),
    June 2024 (no World Cup adjustment).}
    Same layout as Figure~\ref{fig:heatmap}. M1 represents the
    pre-tournament counterfactual: typical June travel volumes
    without the World Cup surge. Dengue $P(\geq 1) > 0.99$ is
    already present at Miami, New York, and Houston under baseline
    conditions, confirming that the World Cup amplifies but does
    not create the dengue importation threat.}
  \label{fig:s1}
\end{figure}

\newpage

% ---- Figure S2 ----------------------------------------------
\refstepcounter{suppfig}
\begin{figure}[htbp]
  \centering
  \includegraphics[width=\textwidth]{figS2_heatmap_wc_IJID.png}
  \caption*{\textbf{Figure S\arabic{suppfig}.}
    \textbf{Importation probability under the WC-adjusted model (M2),
    June 2026.}
    Same layout as Figure~\ref{fig:heatmap}. M2 applies World Cup
    growth factors to national arrivals but distributes all travel
    to US cities via T-100 routing fractions, including the fan
    increment for teams playing in Canada and Mexico. M2 values
    therefore represent a slight overestimate relative to M3 for
    US cities.}
  \label{fig:s2}
\end{figure}

\newpage

% ---- Figure S3 ----------------------------------------------
\refstepcounter{suppfig}
\begin{figure}[htbp]
  \centering
  \includegraphics[width=0.85\textwidth]{figS3_ci_asymmetry_IJID.png}
  \caption*{\textbf{Figure S\arabic{suppfig}.}
    \textbf{Direction of Monte Carlo uncertainty by disease.}
    Each point represents one disease--city combination.
    X-axis: downside ratio ($\lambda_\text{median}/\lambda_\text{lo}$),
    measuring how far the lower CI bound lies below the median.
    Y-axis: upside ratio ($\lambda_\text{hi}/\lambda_\text{median}$),
    measuring how far the upper CI bound lies above the median.
    Points above the diagonal (dashed line) have CIs that extend
    further upward --- the central estimate understates the upper
    plausible value (dengue, influenza, malaria). Points below
    the diagonal have CIs that extend further downward --- the
    central estimate is a conservative upper bound (pertussis).
    The asymmetry arises because the reference parameter values
    differ in their position within the literature range: dengue's
    reference $\rho = 0.10$ is near the lower bound (0.06--0.26);
    pertussis's reference $\rho = 0.10$ is the upper bound
    (0.01--0.10).}
  \label{fig:s3}
\end{figure}

\newpage

% ---- Figure S4 ----------------------------------------------
\refstepcounter{suppfig}
\begin{figure}[htbp]
  \centering
  \includegraphics[width=\textwidth]{figS4_venue_map_IJID.png}
  \caption*{\textbf{Figure S\arabic{suppfig}.}
    \textbf{Geographic distribution of the 16 FIFA World Cup 2026
    stadiums across the United States, Canada, and Mexico.}
    Colours indicate host nation. This analysis focuses on the 11
    US venues (blue); Canadian (Toronto, Vancouver) and Mexican
    (Mexico City, Monterrey, Guadalajara) venues are shown for
    context. The East Coast cluster (New York, Philadelphia, Boston)
    is the densest; Kansas City and Seattle are the most geographically
    isolated US venues.}
  \label{fig:s4}
\end{figure}

\newpage

% ---- Figure S5 ----------------------------------------------
\refstepcounter{suppfig}
\begin{figure}[htbp]
  \centering
  \includegraphics[width=\textwidth]{figS5_diaspora_hubseeding_IJID.png}
  \caption*{\textbf{Figure S\arabic{suppfig}.}
    \textbf{Secondary seeding risk index ($\Omega^B$) for diaspora
    hub cities, by disease (Mechanism B).}
    $\Omega^B_{c,v_\text{match},v_\text{home},d}
      = \lambda_{c,v_\text{match},d}\,\cdot\,\kappa_{c,v_\text{home}}$
    is the expected importation intensity at the match venue weighted
    by the co-national diaspora concentration in the hub city.
    Bars show top 8 hub cities per disease ranked by summed
    $\Omega^B$ across all source countries and match venues.
    This index identifies cities that face secondary risk through
    diaspora members returning from World Cup matches, independently
    of whether they host games themselves.}
  \label{fig:s5}
\end{figure}

\newpage

% ---- Figure S6 ----------------------------------------------
\refstepcounter{suppfig}
\begin{figure}[htbp]
  \centering
  \includegraphics[width=\textwidth]{figS6_country_drivers_omegaA_IJID.png}
  \caption*{\textbf{Figure S\arabic{suppfig}.}
    \textbf{Top 10 source-country contributors to diaspora-network
    importation intensity ($\Omega^A$) per disease.}
    $\Omega^A_{c,\cdot,d} = \sum_v \lambda_{c,v,d}\,\kappa_{c,v}$
    summed across all 11 venue cities. Compare with
    Figure~\ref{fig:countries}: countries whose rank rises here
    (relative to raw $\lambda$) have large, concentrated diaspora
    communities in the cities they tend to travel to; countries
    whose rank falls send many travelers but into cities where
    their community is thin. Colors denote world region.}
  \label{fig:s6}
\end{figure}

\newpage

% ---- Figure S7 ----------------------------------------------
\refstepcounter{suppfig}
\begin{figure}[htbp]
  \centering
  \includegraphics[width=\textwidth]{figS8_diaspora_full_IJID.png}
  \caption*{\textbf{Figure S\arabic{suppfig}.}
    \textbf{Diaspora network importation probability ($P^A(\geq 1)$)
    for all five diseases at 11 US host cities.}
    Derived from $\Omega^A_{v,d} = \sum_c \lambda_{c,v,d}\,\kappa_{c,v}$
    where $\kappa_{c,v} = D_{c,v}/\mathrm{FB}_v$ (ACS 2019--2023).
    Dengue and malaria drive the majority of meaningful diaspora
    signal; measles, pertussis, and influenza show near-zero values
    owing to their very low baseline importation intensities.
    The dengue and malaria rows correspond to Panel A of
    Figure~\ref{fig:diaspora} in the main text.}
  \label{fig:s7b}
\end{figure}

\end{document}
